\scanpage{131}
\mspage{118}
les équivalences
\[
\bigl(u^*(C(I'))\subset C(I)\bigr)\Longleftrightarrow
\bigl(u_*(TF)\subset TF'\bigr)\unsure{,}
\tag*{(6.5.\unsure{5|1})}
\]
\[
\bigl(u^*(TC(I')\subset TC(I)\bigr)\Longleftrightarrow
\bigl(u_*(F)\subset F'\bigr).
\tag*{(6.5.2)}
\]
Les conditions équivalentes (6.5.1) ainsi que
(6.5.2) sont \unsure{satisfaites} quand on est
dans l'un des cas 1°) 2°) 3°) de
6.4. c).

\uline{Dém de 6.4.} Les assertions \textcircled{a}, \textcircled{b}
sont $\pm$ tautologiques, par adjonction,
comparer $\overline{\underline{\mathrm{XIII}}}$ \unsure{prop} 2 et \unsure{cor}, page 46.
Pour prouver c), il suffit de traiter le
cas 3°, car \deleted{1° et 2°} $1^\circ\Rightarrow3^\circ$, $2^\circ\Rightarrow3^\circ$.
Mais le calcul de \deleted{$u_*(X)$}
\[
u_*(X)(i')\simeq\underset{\unsure{i\in|\mathrm{Ob}} I_{/i'}}{\underleftarrow{\lim}}X(i)
\]
montre que si les composantes
connexes $(I_{/i'})_\alpha$ de $I_{/\unsure{i'|j'}}$ ont
chacune un objet final $e_\alpha$, alors
\[
u_*(X)(i')\simeq\prod_\alpha X(e_\alpha).
\]
Or $\Psi$ \inserted{étant clos à droite} est stable par produits quelconques,
donc $X\in\Psi\langle I\rangle$ implique que
\deleted{\unsure{tous les|tt les}} les $u_*(X)(i')$ sont dans $\Psi$, donc
$u_*(X)$ est dans $\Psi\langle I'\rangle$, \unsure{qed.}
% archive sheet number: 130

